\SourceLeaf{248}
\Article{373}
\textsc{Aufgabe.} \emph{Ist die Funktion $X$ modulo $p$ in teilerfremde Faktoren $\xi,\xi',\xi''$ usw. zerlegt, so zerlege man $X$ modulo $pp$ in ähnliche Faktoren $\Xi,\Xi',\Xi''$ usw., sodass}
\[
\xi\equiv\Xi,\quad\xi'\equiv\Xi',\quad\xi''\equiv\Xi'',\text{ usw.}\;(\mathrm{mod.}\,p)
\]

\emph{Lösung.} Es sei $X\equiv\xi\psi(\mathrm{mod.}\,p)$ oder gleichbedeutend $X=\xi\psi+p\Sigma$. Setzt man
\[
\Xi=\xi+p\varphi,\quad\Psi=\psi+p\omega
\]
so ist
\[
\Xi\Psi=X-p\Sigma+(\varphi\psi+\xi\omega)p+pp\varphi\omega
\]
Soll also $\Xi\Psi$ nun $\equiv X(\mathrm{mod.}\,pp)$ sein, so muss $\varphi\psi+\xi\omega-\Sigma$ notwendig durch $p$ teilbar sein. Da aber $\psi$ und $\xi$ modulo $p$ teilerfremde Funktionen sind, lassen sich $\varphi$ und $\omega$ so bestimmen, dass diese Bedingung erfüllt ist (\textsection.\ 336), und zwar außerdem so, dass die Grade von $\varphi$ und $\omega$ jeweils um eins kleiner sind als die Grade der Funktionen $\xi,\psi$. Daher ist $X\equiv\Xi\Psi(\mathrm{mod.}\,pp)$. Offenbar lässt sich $\Psi$ auf ähnliche Weise wieder in Faktoren $\Xi'\Omega$ zerlegen, wobei der Faktor $\Xi'$ dann $\equiv\xi'(\mathrm{mod.}\,p)$ ist, und so fort, bis schließlich
\[
X\equiv\Xi\,\Xi'\,\Xi''\text{ usw.}\;(\mathrm{mod.}\,pp).\quad\text{Q. E. Fac.}
\]

\Article{374}
Hieraus lässt sich leicht beweisen, dass die Funktion $X$ auch modulo $p^3,p^4$ usw. in Faktoren zerlegt werden kann. Allgemeiner sei
\[
X\equiv PQ(\mathrm{mod.}\,p^m)\quad\text{oder}\quad X=PQ+p^mR
\]
und die Funktion $P$ sei zu $Q$ modulo $p$ teilerfremd. Setzt man
\[
P'=P+Ap^m,\quad Q'=Q+Bp^m
\]
so ist
\[
P'Q'=X-p^mR+(AQ+BP)p^m+ABp^{2m}
\]
Hieraus folgt für jeden Modul $p^\nu$ (wobei $\nu$ dann $>m$ und $<2m+1$ ist)
\[
P'Q'\equiv X,\quad\text{wenn}\quad R\equiv AQ+BP(\mathrm{mod.}\,p^{\nu-m})
\]

\SourceLeaf{249}
Daraus erkennt man: Hat die Funktion $X$ modulo $p$ keine gleichen Teiler, so lässt sie sich modulo $p^k$ auf ähnliche Weise in Faktoren zerlegen wie modulo $p$. Hat aber $X$ gleiche Teiler, so wird die Sache viel komplizierter und lässt sich durch die vorigen Prinzipien nicht einmal vollständig behandeln. Da wir nicht alles hierzu Gehörige mitteilen können, wollen wir nur einen einzigen Fall betrachten, der sehr häufig vorkommt und dessen Auflösung nötig ist, um gewisse Zweifel im Vorigen zu beseitigen. Dies ist der Fall, dass nur gleiche Faktoren ersten Grades betrachtet werden. Er lässt sich auch unmittelbar zur Bestimmung der Wurzeln von Kongruenzen verwenden. Allgemein werden wir diesen Gegenstand bei anderer Gelegenheit behandeln.

\Article{375}
Es sei also $X\equiv X'(x-a)^m(\mathrm{mod.}\,p)$, und die Funktion $X'$ sei zu \mbox{$x-a$} teilerfremd: Gesucht sind alle Teiler ersten Grades, die $x-a$ modulo $p$ kongruent sind und $X$ modulo $pp,p^3$ usw. teilen. (Wir setzen voraus, dass die Funktion $X$ ohne einen Modul nicht durch $x-a$ teilbar ist; andernfalls würde $x-a$ die Funktion $X$ nach jedem Modul teilen.) Setzt man $z+a$ für $x$ ein, so erhält man
\[
Z\equiv Z'z^m(\mathrm{mod.}\,p)\quad\text{oder}\quad Z=Z'z^m+pA
\]
Ist nun $Z$ modulo $pp$ durch einen Teiler der Form $z+\alpha p$ teilbar, so muss $A$ notwendig die Form $zZ''+pB$ haben. Ist dies nicht der Fall, so ist die Untersuchung bereits beendet. Setzen wir also
\[
\begin{aligned}
Z&\equiv Z'z^m+pZ''z(\mathrm{mod.}\,pp)\\
&\text{oder}\quad Z=Z'z^m+pZ''z+ppB
\end{aligned}
\]
so ist offenbar $Z$ durch $z$ und jeden anderen diesem modulo $p$ kongruenten Teiler teilbar.

Um die Aufmerksamkeit auf einen bestimmten Fall zu richten, setzen wir $m=4$; man sieht leicht, dass jeder andere Fall auf ähnliche Weise behandelt werden kann. Ist nun $Z$ modulo $p^3$ durch einen Teiler der Form $z+\alpha p$ teilbar, so ist
\[
\begin{aligned}
0&\equiv-\alpha ppZ''+ppB(\mathrm{modd.}\,z+\alpha p,p^3)\\
&\text{oder}\quad\alpha Z''\equiv B(\mathrm{modd.}\,z,p)
\end{aligned}
\]
Nun können drei Fälle eintreten

1) Ist $Z''\equiv0(\mathrm{modd.}\,z,p)$ und $B$ nicht $\equiv0$, so erfüllt offenbar kein Wert von $\alpha$ die Kongruenz, sodass $Z$ modulo $p^3$ keinen Teiler der Form $z+\alpha p$ hat. Die Untersuchung ist somit beendet

\SourceLeaf{250}
2) Ist weder $Z''$ noch $B$ $\equiv0(\mathrm{modd.}\,z,p)$, so hat $\alpha$ einen einzigen Wert, nämlich
\[
\alpha\equiv\frac B{Z''}(\mathrm{modd.}\,z,p).
\]
Daher gibt es einen einzigen Teiler $\equiv z+\alpha p(\mathrm{mod.}\,pp)$ von $Z$ modulo $p^3$, und es ist
\[
Z\equiv V(z+\alpha p)+p^3W
\]
Setzt man nun als Teiler von $Z(\mathrm{mod.}\,p^4)$ den Ausdruck $z+\alpha p+\beta pp$ an, so ist
\[
0\equiv\qquad\qquad
\]
\BookRule{.28\linewidth}
